\documentclass[11pt]{article}

\usepackage[margin=1in]{geometry}
\usepackage{amsmath,amssymb,amsthm}
\usepackage{parskip}
\usepackage{hyperref}

\newtheorem{proposition}{Proposition}
\newtheorem*{remark}{Remark}

\newcommand{\R}{\mathbb{R}}
\newcommand{\E}{\mathbb{E}}
\newcommand{\tr}{\mathrm{tr}}
\newcommand{\diag}{\mathrm{diag}}

\title{A Mathematical Note on SimDINO in Euclidean Embedding Space}
\author{}
\date{}

\begin{document}
\maketitle

\begin{abstract}
This note analyzes the SimDINO-style objective used in the mini experiment as a regression-based self-distillation method in Euclidean space with an explicit spectral anti-collapse regularizer. The central claim is that the objective is mathematically cleaner than the other losses in the repository: it combines Euclidean alignment, per-dimension scale control, and a log-determinant correlation-rate term that spreads information across all directions of the representation. The resulting geometry is best understood as a whitened Euclidean latent space rather than a probabilistic Gaussian embedding model.
\end{abstract}

\section{Setup}

Let the student produce two view embeddings $s_1, s_2 \in \R^D$ and the teacher produce $t_1, t_2 \in \R^D$. Write the batch of student features as
\begin{equation}
    Z \in \R^{N \times D},
\end{equation}
where $N$ is the number of student embeddings entering the regularizer. Define the batch mean and per-dimension standard deviation by
\begin{equation}
    \mu = \frac{1}{N} \sum_{n=1}^{N} Z_n,
    \qquad
    \sigma_i^2 = \frac{1}{N} \sum_{n=1}^{N} (Z_{ni} - \mu_i)^2,
\end{equation}
and the whitened batch
\begin{equation}
    \widehat{Z} = (Z - \mu)\diag(\sigma)^{-1}.
\end{equation}
The empirical correlation matrix is then
\begin{equation}
    C = \frac{1}{N}\widehat{Z}^{\top}\widehat{Z} \in \R^{D \times D}.
\end{equation}

The intended SimDINO objective can be written as
\begin{equation}
    \mathcal{L}
    =
    \mathcal{L}_{\mathrm{align}}
    +
    \lambda_{\mathrm{sig}} \mathcal{L}_{\mathrm{sig}}
    -
    \gamma R_{\mathrm{corr}},
\end{equation}
where
\begin{align}
    \mathcal{L}_{\mathrm{align}}
    &=
    \frac{1}{2}
    \left(
    \|s_1 - t_2\|_2^2 + \|s_2 - t_1\|_2^2
    \right), \\
    \mathcal{L}_{\mathrm{sig}}
    &=
    \frac{1}{D}
    \sum_{i=1}^{D} (\sigma_i - 1)^2, \\
    R_{\mathrm{corr}}
    &=
    \frac{1}{2}
    \log \det(I + \beta C),
    \qquad
    \beta = \frac{D}{\varepsilon^2}.
\end{align}

\section{Spectral Interpretation}

The key term is the correlation-rate objective. Since $C$ is symmetric positive semidefinite, let
\begin{equation}
    C = Q \diag(\lambda_1, \dots, \lambda_D) Q^{\top},
    \qquad
    \lambda_i \ge 0.
\end{equation}
Then
\begin{equation}
    R_{\mathrm{corr}}
    =
    \frac{1}{2}
    \sum_{i=1}^{D} \log(1 + \beta \lambda_i).
\end{equation}

Because each dimension of $\widehat{Z}$ has unit empirical variance, the trace is approximately fixed:
\begin{equation}
    \tr(C)
    =
    \sum_{i=1}^{D} \lambda_i
    \approx D.
\end{equation}
Therefore the optimization problem induced by $R_{\mathrm{corr}}$ is: maximize a concave function of the eigenvalues subject to a fixed sum and nonnegativity constraints.

\begin{proposition}
Among all positive semidefinite matrices $C$ with fixed trace $\tr(C)=D$, the quantity
\begin{equation}
    \sum_{i=1}^{D} \log(1 + \beta \lambda_i)
\end{equation}
is maximized when
\begin{equation}
    \lambda_1 = \cdots = \lambda_D = 1,
\end{equation}
equivalently $C = I$.
\end{proposition}

\begin{proof}
The function $f(x)=\log(1+\beta x)$ is strictly concave on $x \ge 0$. By Jensen's inequality,
\begin{equation}
    \frac{1}{D}\sum_{i=1}^{D} f(\lambda_i)
    \le
    f\left(
        \frac{1}{D}\sum_{i=1}^{D}\lambda_i
    \right)
    =
    f(1),
\end{equation}
with equality if and only if all $\lambda_i$ are equal. Since the trace constraint gives $\frac{1}{D}\sum_i \lambda_i = 1$, the maximum is attained at $\lambda_i = 1$ for all $i$.
\end{proof}

So the regularizer does more than suppress pairwise correlations. It pushes the full spectrum toward isotropy and therefore penalizes rank collapse directly.

\section{Why the Log-Det Term is Strong}

Simple decorrelation losses often penalize only off-diagonal entries of $C$. By contrast, the log-det term acts on the entire spectrum. Its derivative with respect to an eigenvalue is
\begin{equation}
    \frac{\partial R_{\mathrm{corr}}}{\partial \lambda_i}
    =
    \frac{1}{2}\frac{\beta}{1 + \beta \lambda_i}.
\end{equation}
This derivative is larger for smaller $\lambda_i$, which means the regularizer pushes especially hard on underrepresented directions. That is why it behaves as an anti-collapse mechanism rather than merely a redundancy penalty.

\section{Asymptotic Regimes}

The hyperparameter $\varepsilon$ controls the strength of the spectral term through $\beta = D/\varepsilon^2$.

\subsection{Weak-rate regime}
If $\beta$ is small, then
\begin{equation}
    \log \det(I + \beta C)
    =
    \tr(\log(I+\beta C))
    \approx
    \beta \tr(C),
\end{equation}
so
\begin{equation}
    R_{\mathrm{corr}}
    \approx
    \frac{\beta}{2}\tr(C).
\end{equation}
Since $\tr(C)$ is nearly fixed after standardization, the rate term becomes almost constant and stops providing a useful optimization signal.

\subsection{Strong-rate regime}
If $\beta$ is large, then
\begin{equation}
    \log \det(I + \beta C)
    =
    D \log \beta + \log \det(C) + o(1).
\end{equation}
Up to an additive constant, the objective then maximizes $\log \det(C)$, which is the log-volume spanned by the whitened batch. This heavily penalizes small eigenvalues and strongly discourages low-rank structure.

\begin{remark}
This explains why $\varepsilon$ is not a cosmetic hyperparameter. Large $\varepsilon$ weakens the spectral regularizer until it is nearly inert, while small $\varepsilon$ pushes the model toward aggressively isotropic representations.
\end{remark}

\section{Geometric Interpretation}

The SimDINO objective is best read as the combination of three geometric biases:
\begin{enumerate}
    \item Euclidean alignment: $\mathcal{L}_{\mathrm{align}}$ pulls paired views together directly in $\R^D$.
    \item Scale control: $\mathcal{L}_{\mathrm{sig}}$ keeps each coordinate near unit marginal standard deviation.
    \item Spectral spreading: $R_{\mathrm{corr}}$ distributes representation energy across all principal directions.
\end{enumerate}

Taken together, these terms favor a representation that is view-consistent, non-collapsed, and approximately whitened. That is a coherent Euclidean learning objective. In particular, it avoids the simplex geometry induced by softmax-based DINO objectives and instead works directly with continuous coordinates in $\R^D$.

\section{What the Objective is Not}

Although the geometry is Gaussian-like, this is not yet a probabilistic Gaussian embedding model. The method does not predict a per-sample mean and covariance, and it does not match distributions through a divergence such as KL or Wasserstein distance. Instead, it learns deterministic vectors whose batch statistics are pushed toward an isotropic structure.

So the most accurate description is:
\begin{quote}
Euclidean self-distillation with spectral whitening.
\end{quote}
rather than
\begin{quote}
probabilistic Gaussian embeddings.
\end{quote}

\section{Implementation Notes}

The mathematical story above treats the intended objective. In the current mini-experiment implementation there is one detail worth flagging: the alignment term is scaled as
\begin{equation}
    0.5 \cdot \frac{\|s_1 - t_2\|_2^2 + \|s_2 - t_1\|_2^2}{2},
\end{equation}
which is smaller by a factor of $2$ than the natural average of the two cross-view losses. This means the current code weights the regularizer more strongly relative to alignment than the written objective would suggest.

\section{Testable Predictions}

The analysis suggests a few concrete predictions.

\begin{enumerate}
    \item If $\varepsilon$ is too large, performance should degrade because the code-rate term approaches a constant.
    \item If $\varepsilon$ is reduced moderately, the smallest eigenvalues of $C$ should increase and the spectrum should flatten.
    \item The method should be more sensitive to batch size than a plain MSE objective, because $C$ is a batch-level estimate.
    \item Replacing raw MSE alignment with normalized alignment should reduce conflicts between alignment and scale control.
\end{enumerate}

\section{Conclusion}

From a mathematical perspective, SimDINO is the most coherent Euclidean objective in the current repository. Its core regularizer is not just a decorrelation term; it is a spectral objective that explicitly rewards full-rank, isotropic representations under fixed marginal scale. That makes it a principled foundation for learning in flat embedding spaces. The main caution is terminological: the method currently learns whitened deterministic embeddings, not full Gaussian distributions. If a genuinely Gaussian embedding model is desired, the next step is to parameterize sample-wise uncertainty explicitly.

\end{document}
